% content.tex — the body of paper_05_subject_formation, as a book chapter.
% GENERATED by split_content.py. The standalone paper still builds from
% paper_05_subject_formation.tex; both read this file, so they cannot drift apart.

\section{Introduction: The Question of the Intimate Subject}
% =====================================================================

Who is formed in an intimate relationship? The question sounds rhetorical, and most moral and legal theory treats it as settled before theory begins. Contract doctrine presupposes parties; rights discourse presupposes holders; even the richest theories of love tend to presuppose two completed persons who then exchange affection, promises, and goods. The previous papers of this series have already strained against this presupposition. The third paper argued that a vow made in an intimate relationship cannot be unilaterally legislated, because the meaning of the vow's central terms is held in common and must remain answerable to the beloved.\footnote{Paper III of this series, \emph{The Normativity of Self-Legislation in Intimate Relationships: Freedom, Morality, and Justice} (manuscript, 2026).} The fourth paper argued that the gift in intimate life cannot be settled or cleared like a debt, because its value is produced relationally, by its position in a circulation that no single party owns.\footnote{Paper IV of this series, on the gift and the circulation of value in intimate life (manuscript, 2026).} Both arguments quietly presuppose a stronger claim that neither paper defended: that the \emph{subjects} of intimate life, the very ones who vow and give, are themselves relationally constituted. This paper defends that claim, names the mechanism, and traces its consequences for ethics, for the structure of rights, and for the theory of justice.

The thesis is this:

\begin{prop}[Mechanism Thesis]\label{p05:prop:mechanism}
The subject of intimate life is formed and continuously sustained through the presence of the other and the practice of joint attention: the recursive, mutually open structure in which each party attends to a shared object and to the other's attending. The other and joint attention are not influences upon an already-formed subject; they are the mechanism of its formation.
\end{prop}

Four consequences follow, and they occupy the second half of the paper. First, if joint attention is the mechanism of subject formation, then the contemporary industrial organization of attention, what is loosely called the attention economy, is not merely a nuisance of distraction. It competes with intimate life \emph{at the level of constitution}, extracting precisely the material through which subjects form one another, and the competition is rigged in its favor; a regenerative criterion of relational sustainability falls out of the analysis. Second, the symbolic order enters the account a second time, no longer only as a mechanism of formation but as the standing \emph{distributor} of attention: every adult triangle contains a virtual fourth vertex that pre-allocates salience, and this yields the concepts of attentional alienation, attention executed under a script one cannot appropriate, and attentional autonomy, the joint re-authorship of that script, which proves to be the third paper's self-legislation continued by attentional means. Third, the question of what each partner owes the other in attention acquires a precise jurisprudential form: is there a claim-right to attention? The answer defended here is no, for reasons isomorphic to the third paper's prohibition on unilateral legislation, and the positive account that replaces the claim-right is an imperfect duty in Kant's register, a virtue in Murdoch's, and a structural criterion of attentional justice with a distributional clause. Fourth, the theory of justice acquires a category it has lacked: \emph{constitutive injustice}, the wrong done not to a person's holdings, opportunities, or enumerated rights, but to the conditions under which that person becomes and remains a subject at all.

\subsection{Method and the no-overclaim discipline}

The series treats the intimate dyad as the minimal nontrivial substructure of the social: the model organism of relational theory. But the discipline announced in the earlier papers binds here with special force, because this paper, alone in the series so far, rests part of its case on empirical developmental psychology. Findings about infants do not transfer to adult lovers by default. Section \ref{p05:sec:transfer} therefore argues the transfer explicitly, identifies what does not carry over (with Simmel's analysis of the dyad as the principal source of resistance), and states the residual claim at its true strength. Where the paper coins a concept, it marks the coinage; where it relies on a contested empirical literature, it flags the contest; and where its argument is conditional, the condition is stated in the proposition itself.

Two boundaries are drawn in advance, so that their observance later is read as discipline rather than omission. First, the paper treats the dyad alone: the entrance of third \emph{persons} into the triangle, the child who is at once shared object and emergent subject, the rival, transforms the structure in ways Simmel's analysis of the triad already predicts \parencite{simmel1950}, and is deferred to later work in the series. Second, the embodied and erotic dimensions of intimate attention are set aside: not as unimportant, but because they demand a phenomenological treatment, and an experiential warrant, that this paper does not claim to possess; the series' standing rule is to argue only from cases whose stakes the author can answer for.

A final methodological note concerns vocabulary. The framework of this series has been developed under the name \emph{relational being}, and that name is shared, coincidentally, with the title of Kenneth Gergen's 2009 book \parencite{gergen2009}. The overlap is more than nominal: Gergen too denies the bounded, self-contained self and relocates the origin of meaning in relational process. The differences are stated in §\ref{p05:sec:gergen}; the present use of the term is independent of and not derived from Gergen's, but his priority in the title is acknowledged, and the engagement is substantive rather than defensive.

% =====================================================================
\section{Theories of the Relational Subject: Three Types of Constitutive Mechanism}
\label{p05:sec:survey}
% =====================================================================

A survey organized by school would set phenomenology beside psychoanalysis beside developmental psychology and let the reader draw the lines. This section is organized differently, by the \emph{type of mechanism} each account proposes for the constitution of the subject, because the paper's central argument turns on a structural relation among the mechanisms themselves. Three types recur across otherwise distant literatures: mirror mechanisms, in which the subject is formed in the reflecting regard of an other; triadic mechanisms, in which the subject is formed by sharing a world with an other; and symbolic mechanisms, in which the subject is formed by induction into a system of signs that precedes it. The claim to be defended in Section \ref{p05:sec:thesis} is that the triadic type is the hinge between the other two, and that accounts which omit it cannot explain how a creature formed in mirrors comes to inhabit a language.

\subsection{Mirror mechanisms: the subject formed in the regard of the other}
\label{p05:sec:mirror}

The oldest modern statement is sociological. Cooley's looking-glass self holds that self-feeling arises from the imagination of our appearance to the other and of the other's judgment of that appearance \parencite{cooley1902}. Mead gave the idea its enduring structure: the self divides into the ``I'' that acts and the ``me'' that is the internalized attitude of others, so that self-consciousness is, constitutively, taking the perspective of the other upon oneself \parencite{mead1934}. In both accounts the mechanism is specular: the other functions as a surface in which a self, otherwise invisible to itself, first appears.

Psychoanalysis radicalized the figure. Lacan's mirror stage places the formation of the \emph{I} in the infant's jubilant identification with its specular image: an anticipation of bodily unity that the infant's actual motor incapacity belies \parencite{lacan2006}. Two features matter for everything that follows. First, the identification is a \emph{m\'econnaissance}, a misrecognition: the unity the infant assumes is not yet, and in a sense never will be, its own. The subject is founded on an image it mistakes for itself, and alienation is therefore not an accident that befalls the subject but a structural moment of its formation. Second, the mirror need not be glass. Winnicott, correcting and absorbing Lacan, located the precursor of the mirror in the mother's face: the infant looks at the mother and sees, in her expression, itself being seen \parencite{winnicott1971}. His earlier, more famous remark gives the ontological version: there is no such thing as an infant, meaning that wherever one finds an infant one finds maternal care, a nursing couple rather than a unit \parencite{winnicott1960}. The minimal subject is already a dyad.

The mirror mechanisms share a structure: they are \emph{dyadic} and \emph{reflective}. Two terms face one another, and constitution runs through the image of the one in the regard of the other. What the mirror accounts lack, and know they lack, is the world. Nothing in the specular relation explains how the two come to be jointly directed at anything beyond the relation; the mirror is a closed circuit. Lacan's own solution was to break the circuit from above, by the symbolic order (§\ref{p05:sec:symbolic}). The literature surveyed next breaks it from below.

\subsection{Triadic mechanisms: the subject formed by sharing a world}
\label{p05:sec:triadic}

In 1975 Scaife and Bruner reported that infants follow the direction of an adult's gaze, and dated the emergence of the capacity within the first year \parencite{scaife1975}. The finding seeded a research program whose mature shape is the study of \emph{joint attention}: not two creatures looking at the same thing, which can happen by accident, but two creatures attending to the same thing \emph{together}, each registering the other's attention, in a structure that is mutually open \parencite{tomasello1995}. Developmental psychology marks a transition, often called the nine-month revolution, in which the infant moves from dyadic engagement with persons (protoconversation, mutual gaze, affect exchange: what Trevarthen called primary intersubjectivity) to triadic engagement in which infant, adult, and object form a referential triangle: the infant looks where the adult looks, alternates gaze between object and adult's eyes, points, shows, and checks back \parencite{trevarthen1979, trevarthen1978}. Trevarthen and Hubley named the achievement secondary intersubjectivity; Tomasello and colleagues argue it rests on a species-distinctive capacity for \emph{shared intentionality}: the ability and motivation to form, with another, a ``we'' that attends and intends as one \parencite{tomasello2005, tomasello2014}. Bruner showed that these triangles are the cradle of language: reference is learned inside formats of joint attention, the give-and-take of looking, naming, and checking, before any word is understood \parencite{bruner1983}. And the dyadic precursor is itself robustly interactive rather than merely specular: in the double-video experiments of Murray and Trevarthen, two-month-olds presented with a replayed (and thus non-contingent) image of their own mother became distressed and disengaged, although the perceptual stimulus was identical \parencite{murray1985}. What the infant seeks is not the sight of the mother but her \emph{live responsiveness}: presence, not image.

Philosophy reached the same structure by an independent route. Davidson's triangulation argument holds that thought with objective content, thought that can be \emph{about} something and can be \emph{wrong} about it, requires a triangle: two creatures interacting with each other and with a shared world, each line of the triangle indispensable for fixing what any thought is about \parencite{davidson2001}. The argument is transcendental rather than developmental: it does not describe how children in fact acquire objectivity but argues that nothing could count as objective thought outside the triangular structure. Its convergence with the developmental findings is therefore significant. The psychology says: this is how subjects are in fact formed. Davidson says: there is no other way they could be. Together they elevate joint attention from an empirical milestone to a candidate \emph{constitutive condition} of subjectivity, which is precisely the rank the Mechanism Thesis assigns it.

The triadic mechanisms differ from the mirror mechanisms in two respects that the argument of Section \ref{p05:sec:thesis} will exploit. First, they are \emph{world-involving}: the third vertex of the triangle is not another regard but an object, and constitution runs through co-orientation toward it. Second, they are \emph{veridical} where the mirror is illusory: joint attention, when achieved, is really achieved; there is no structural m\'econnaissance in pointing at the moon together. The tension between this veridicality and Lacan's constitutive misrecognition is real, and §\ref{p05:sec:lacantomasello} is devoted to it.

\subsection{Symbolic mechanisms: the subject formed by the sign}
\label{p05:sec:symbolic}

The third family locates constitution in language and the social systems it carries. Saussure's structural insight, that the value of a sign is purely differential, fixed by its position in a system rather than by any intrinsic bond to its object, dissolved the picture of language as a nomenclature adopted by ready-made minds \parencite{saussure1959}. If meaning is systemic, the speaking subject is an effect of its position in the system, not the system's author. Lacan translated the insight into psychoanalysis: the symbolic order, the big Other, precedes the subject, assigns it a name, a kinship position, a desire mediated by the desire of the Other; the subject that emerges is divided by the very signifiers that represent it \parencite{lacan2006}. Althusser supplied the political register: ideology interpellates individuals as subjects, the policeman's hail on the street that one turns to answer being the elementary scene in which subjection and subjectivation coincide \parencite{althusser1971}. Butler radicalized the temporality of the scene: the subject is not interpellated once but must be reiterated, performatively, and the necessity of reiteration is also the possibility of subversion \parencite{butler1997}.

Within developmental psychology the symbolic family has a precise representative, and he is the natural bridge to the triadic findings. Vygotsky's genetic law of cultural development states that every higher psychological function appears twice: first between people, as an interpsychological category, and only then within the child, as an intrapsychological category \parencite{vygotsky1978}. Inner speech is internalized dialogue; private thought is a social relation folded inward. The law is the most direct empirical formulation in the literature of the claim that the other is the mechanism, not the occasion, of subject formation, and it shows the symbolic mechanism to be developmentally \emph{downstream}: what is internalized must first have been shared, and what is shared must first have been jointly attended.

\subsection{Theory of mind and the enactivist correction; Gergen distinguished}
\label{p05:sec:gergen}

One influential research program appears to compete with the account being assembled, and must be addressed rather than enlisted. The theory-of-mind paradigm, descending from Premack and Woodruff's question about the chimpanzee \parencite{premack1978} and consolidated by the false-belief studies \parencite{baroncohen1985}, models social understanding as the attribution of hidden mental states to others: the child as a small theorist, inferring invisible beliefs behind observed behavior. On this picture the other is an \emph{epistemic problem} for an already-constituted subject, and intersubjectivity is mind\emph{reading}: third-personal, spectatorial, inferential.

The enactivist critique inverts the picture. Gallagher argues that in ordinary second-person interaction we do not infer mental states but \emph{perceive} intention and emotion directly in expressive behavior, and that the theorizing stance is a late, specialized, and comparatively rare achievement \parencite{gallagher2008}. De Jaegher and Di Paolo's participatory sense-making generalizes the point: social understanding is constituted in the coordination dynamics of interaction itself; meaning is not transmitted between two closed minds but generated between them, in a process to which both contribute and which neither controls \parencite{dejaegher2007}. Reddy's developmental work supplies the second-person evidence: infants engage with attention directed \emph{at them} months before they pass any test of attributing attention to others, suggesting that the lived, second-person relation is prior to the third-person theory of it \parencite{reddy2008}. For the present argument the critique is a resource, not a difficulty: it relocates the formation of social understanding exactly where the Mechanism Thesis needs it, in the interaction, and it demotes the spectatorial stance that would make the other an object of inference rather than a partner in constitution. Theory of mind remains in the account as a real but derivative competence, a tool subjects use, not the process that makes them subjects.

Finally, the nominal collision noted in the introduction. Gergen's \emph{Relational Being} argues that bounded selfhood is a damaging cultural construction and that all meaning, including the meaning of selfhood, originates in relational process \parencite{gergen2009}. The present framework shares the negative thesis and honors the convergence. It differs in three commitments. First, Gergen's program is social-constructionist all the way down and is correspondingly reluctant to specify mechanisms; this series is mechanist, and the present paper's entire purpose is to name one. Second, Gergen's unit is the diffuse field of relational process; the present unit is the structured triangle, with determinate vertices and a determinate recursion. Third, this series binds its ontology to a theory of justice with jurisprudential ambitions (Hohfeldian analysis, criteria, the architecture of rights), where Gergen's normative register is therapeutic and dialogic. The shared name marks a shared enemy, the self-contained self, rather than a shared theory.

\subsection{The map, assembled}

The three mechanism-types now stand in a suggestive arrangement. Mirror mechanisms are dyadic and pre-linguistic; symbolic mechanisms are linguistic and social at a scale far beyond the dyad; and the triadic mechanisms sit exactly between, sharing the face-to-face intimacy of the mirror and the world-directedness of the sign. Developmentally they occupy the interval between the mirror stage and language. Structurally, the shared object of the triangle is the first entity to stand outside the closed circuit of mutual regard, and the formats of joint attention are the scaffolding within which signs are first acquired. The next section argues that this arrangement is not an artifact of exposition but the architecture of subject formation itself.

% =====================================================================
\section{The Thesis: The Other and Joint Attention as the Mechanism of Formation}
\label{p05:sec:thesis}
% =====================================================================

\subsection{The triadic as the missing middle term}

Consider the explanatory gap each family leaves open. The mirror accounts end with a subject captivated by an image: dyadic, closed, worldless. The symbolic accounts begin with a subject already inside language: addressed, named, positioned. Between the captivated infant and the interpellated speaker lies a transition that neither family explains. How does a creature formed in mirrors come to be formable by signs? Interpellation presupposes that the hail can be \emph{recognized as addressed to me about something}; the mirror relation contains no resources for aboutness at all.

The triadic mechanism is the answer, and the developmental record is unusually kind to the philosophical point. The referential triangle is the first structure in which a second person and a world co-occur: the infant attends to the object, attends to the adult's attending, and checks the alignment. Reference, the elementary capacity presupposed by every symbolic mechanism, is assembled inside these triangles \parencite{bruner1983}; Davidson's argument adds that it could not be assembled anywhere else \parencite{davidson2001}. Vygotsky's law then describes the direction of travel: the interpersonal triangle is folded inward and becomes the intrapersonal structure of thought \parencite{vygotsky1978}.

\begin{prop}[Hinge Thesis]\label{p05:prop:hinge}
The triadic mechanism is the developmental and structural hinge between mirror and symbolic constitution: it opens the dyadic mirror onto a world, and it builds the scaffolding within which symbolic induction becomes possible. Accounts of relational subject formation that omit it are not merely incomplete; they are disconnected, a beginning and an end with no middle.
\end{prop}

Two clarifications guard the thesis. First, ``hinge'' is not ``replacement'': the claim is not that joint attention supersedes mirror or symbol, but that all three mechanisms remain operative in the formed subject, in a layered structure to be described presently. Second, the thesis is stated for human ontogeny and defended at that level; its extension to adult intimacy is a separate argument, made under separate discipline in Section \ref{p05:sec:transfer}.

For the purposes of later sections it will help to fix the structure formally, in the minimal notation the series prefers. Write $\mathrm{Att}(a,x)$ for ``$a$ attends to $x$.''

\begin{defn}[Joint attention]\label{p05:def:ja}
Parties $a$ and $b$ are in joint attention on object $x$, written $\mathrm{JA}(a,b,x)$, iff
(i) $\mathrm{Att}(a,x)$ and $\mathrm{Att}(b,x)$;
(ii) $\mathrm{Att}(a,\mathrm{Att}(b,x))$ and $\mathrm{Att}(b,\mathrm{Att}(a,x))$;
(iii) the structure in (i)--(ii) is mutually open: neither party's awareness of it is hidden from the other, and each may, at any level, check and find the alignment confirmed.
\end{defn}

Clause (iii) deliberately avoids requiring an actual infinite hierarchy of metarepresentations; what it requires is \emph{openness}, the standing availability of the next check, which is how the developmental literature itself characterizes the phenomenon \parencite{tomasello1995}. Three features of the definition do later work. Joint attention is \emph{relational} in the strict sense: no conjunction of facts about $a$ alone and $b$ alone entails $\mathrm{JA}(a,b,x)$, because clause (iii) is a fact about the between. It is \emph{non-delegable}: $a$ cannot satisfy her side by proxy, since what (ii) requires is her attending, not an attending procured by her. And it is \emph{fragile}: it can fail silently, with both parties attending to $x$ while the mutual openness has lapsed, which is why presence and its simulacra (Section \ref{p05:sec:parasite}) can come apart.

\subsection{Reconciling Lacan and Tomasello: a layered account}
\label{p05:sec:lacantomasello}

The survey left a standing tension. Lacan's mirror founds the subject on misrecognition; the joint-attention literature describes real, achieved, mutually confirmed alignment. One tradition makes alienation structural; the other makes cooperation primordial. A paper that enlisted both without comment would deserve the referee's suspicion, so the reconciliation is made explicit.

The proposal is a \emph{layered} account. The imaginary layer (mirror) and the triadic layer (joint attention) are distinct, developmentally sequenced, and never abolished by what succeeds them; the symbolic layer is built upon the triadic and reorganizes both. In the formed subject all three run concurrently. The layers are individuated by their objects and their failure modes. The imaginary layer takes as object the \emph{image} of self and other, and its characteristic distortion is captation: the lure of the unified image, the projection onto the other of what completes me. The triadic layer takes as object the \emph{shared world}, and its characteristic achievement is veridical co-orientation; its failure is not illusion but lapse, the silent collapse of mutual openness. The symbolic layer takes as object the \emph{position}: name, role, status in the differential system; its characteristic pathology is the subject's reduction to its position.

On this account there is no contradiction in asserting both that the lover's gaze is shot through with imaginary projection and that the lovers' attention to a shared world is really shared. The two claims are made about different layers, and adult intimacy is precisely the condition in which both run at maximal intensity, simultaneously. When I look at her, the imaginary layer is active: Lacan's dictum that ``you never look at me from the place from which I see you'' \parencite{lacan1998} names a gap that no intimacy closes, and the fantasy that it could be closed is itself the imaginary at work. But when we look together at the same moon, the triadic layer is active, and its achievement is not undone by the imaginary gap: the moon is genuinely shared even though the gaze is not symmetrical. The mature ideal this paper will defend for intimate presence (§\ref{p05:sec:presence}) is built exactly here: not the fantasy of full mutual transparency, which the imaginary layer forbids, but veridical co-orientation \emph{sustained in the acknowledged presence of the gap}.

The layering also disciplines each literature with the other. Against an uncritical Tomasellian optimism, the imaginary layer insists that shared intentionality in adults is never innocent of projection: the ``we'' that intends together is also, always, a ``we'' each party images in his own style. Against a totalized Lacanian pessimism, the triadic layer insists that misrecognition is not the whole of relation: there are achievements of alignment that are not lures, and a theory in which nothing is ever really shared cannot explain the most ordinary scenes of human learning, the picture book, the pointed finger, the checked glance.

\subsection{The first third: a Levinasian difficulty answered}
\label{p05:sec:third}

Levinas poses the sharpest external challenge to any justice-theoretic treatment of the dyad. For Levinas, the face-to-face relation is ethically primary and radically asymmetrical: my responsibility for the other is infinite, non-reciprocal, and prior to any contract \parencite{levinas1969}. Justice, by contrast, requires comparison, measure, symmetry, and these enter, Levinas says, only with the \emph{third party}, le tiers: when there are three, I must weigh one other against another, and the incomparable must be compared \parencite{levinas1998}. The difficulty for this series is immediate: its entire program treats the dyad as a site of justice, of rights structure, of fair and unfair. If justice begins with the third, and the dyad has no third, then either the program is confused or the dyad must contain a third after all.

The triadic analysis supplies the answer, and it is not a trick. The dyad of joint attention is never merely a dyad: its constitutive structure contains a third vertex, the shared object. The first third party is the world. This is not yet Levinas's tiers, a third \emph{person} with a face of his own, and the difference is respected: the shared object makes no claims and suffers no wrongs. But it performs the structural function Levinas assigns the third: it breaks the closed circuit of the face-to-face, institutes a common term to which both parties stand in comparable relations, and thereby makes alignment, and so misalignment, and so \emph{complaint}, first possible. Whose concern the shared gaze serves, whose project the common object belongs to, whose world the two are jointly inhabiting: these are questions of justice, and they are well-formed inside the dyad because the dyad's own formative mechanism is triangular. When, in Section \ref{p05:sec:transfer}, the relation itself becomes the shared object, the point sharpens further: the ``we'' under joint attention is a quasi-third with a grammar of claims of its own.

The reconciliation preserves what is right in Levinas, the asymmetry of the ethical, which survives in a register of its own that the triangle neither generates nor cancels, while denying the inference that the dyad is a justice-free zone. The face exceeds measure; the triangle institutes it; intimate life is the standing coexistence of the two.

% =====================================================================
\section{Transfer to Adult Intimacy: The Reflexive Triangle}
\label{p05:sec:transfer}
% =====================================================================

\subsection{What does not transfer}

The joint-attention literature is a literature about infants, and the no-overclaim discipline forbids silent extrapolation. Three disanalogies are acknowledged at full strength. First, asymmetry of formation: the infant is being formed, the caregiver already formed; in adult intimacy both parties arrive already constituted, with stabilized imaginaries and full symbolic positions. Whatever formation occurs between adults is re-formation, maintenance, and gradual re-weighting, not origination. Second, the dyad's structural peculiarity: Simmel showed that the dyad differs in kind from all larger groups, possessing no supra-individual structure that survives the exit of either member, no majority, no coalition, no delegation of the relation itself \parencite{simmel1950}. Claims developed for the infant-caregiver pair, which is embedded in a family and a society that outlast it, do not automatically hold for a structure whose existence is identical with the participation of exactly two. Third, the mediating role of language: adult joint attention is saturated with the symbolic in a way infant triangles are not; adults share objects under descriptions, and descriptions are contestable in ways gaze is not.

The transferred claim is therefore stated conditionally and at reduced strength:

\begin{prop}[Transfer Thesis]\label{p05:prop:transfer}
In adult intimacy, joint attention is not the origin of the subject but the mechanism of its ongoing maintenance and renegotiation: the process by which each partner's self-understanding, standing, and horizon of concern are continuously re-formed in the presence of the other. The mechanism operates at lower intensity and higher mediation than in ontogeny, but it does not change in kind: Definition \ref{p05:def:ja} applies unchanged, and its three features (relationality, non-delegability, fragility) carry over fully.
\end{prop}

The defense is that the alternative is incredible. To deny the Transfer Thesis is to hold that the formed subject is finished: that nothing in the adult's selfhood remains answerable to the quality of attention in her closest relations. Everything ordinary testifies against this: the observable expansion of a person who is, over years, genuinely seen, and the equally observable contraction of one who is not; the way long partners co-author each other's memory, taste, and idiom; the clinical commonplace that the withdrawal of attention in a marriage is experienced not as a loss of a good but as a loss of \emph{reality}. The double-video finding gives the deep reason \parencite{murray1985}: what the human seeks in the other is live responsiveness, and there is no age at which this need is outgrown, only ages at which its frustration is better disguised.

\subsection{The relation itself as shared object: the reflexive triangle}
\label{p05:sec:reflexive}

The original contribution this section offers is a thesis about the \emph{mature form} of intimate joint attention. In ontogeny the triangle is (self, other, thing): the picture book, the moon. In adult intimacy a further structure becomes available, and its availability is, the paper proposes, the criterion of the relationship's maturity:

\begin{defn}[Reflexive triangle]\label{p05:def:reflexive}
A reflexive triangle obtains between $a$ and $b$ iff $\mathrm{JA}(a,b,R_{ab})$, where $R_{ab}$ is the relation between $a$ and $b$ itself: the partners jointly attend to the ``we,'' each attending to it and to the other's attending, in mutual openness.
\end{defn}

A regress worry arises immediately: $R_{ab}$ includes the parties' attentional dispositions, so joint attention to $R_{ab}$ is attention to a structure that contains that very attention. The worry is dissolved by stratification. Let $R^{0}_{ab}$ be the relation as constituted prior to a given reflexive act, and let each act of joint attention to $R^{n}_{ab}$ contribute to $R^{n+1}_{ab}$. Reflexive attention is then always attention to the relation \emph{as it stands}, and its own occurrence becomes part of what stands for the next act. There is no vicious circle; there is a spiral, and the spiral is the phenomenon: relationships that attend to themselves change by that attending, which is why the practice is formative rather than merely descriptive.

The reflexive triangle is the point at which the series' earlier results converge, and the convergence is evidence for the analysis. The vow of Paper III is, in the present vocabulary, a device for placing the relation under joint attention: it makes $R_{ab}$ an explicit object, before both parties, in mutual openness, and the third paper's conclusion, that the vow's terms must remain answerable to the beloved, is clause (iii) of Definition \ref{p05:def:ja} in normative dress. The gift of Paper IV is likewise a joint-attention device: its meaning function was shown there to be non-factorizable across the parties, and the genetic explanation is now available: the gift means what it means only inside the triangle, as an object the relation attends to together, which is why its value cannot be settled bilaterally like a price. Anniversaries, rituals, shared tables, the deliberate review of a quarrel, the conversation that begins ``what are we becoming'': all are technologies of the reflexive triangle. A relationship without them may persist, but it persists as habit; it is no longer forming its subjects, only housing them.

\subsection{The presence of the other: attentional presence and the acknowledged gap}
\label{p05:sec:presence}

The phrase ``the presence of the other'' in this paper's title can now be made precise, and distinguished from two counterfeits. Physical co-presence is not presence: two bodies in one room, each absorbed elsewhere, satisfy no clause of Definition \ref{p05:def:ja}. Surveillance is not presence either: one-directional monitoring satisfies clause (i) for one party while violating the mutual openness of clause (iii) entirely; the watched party's attending is captured, not joined. Presence in the constitutive sense is \emph{attentional presence}: availability for the triangle, the standing offer of clauses (i) through (iii).

Buber's distinction between the I-Thou and I-It relations marks the same boundary from within dialogical philosophy: the Thou is not an object of experience but a partner in relation, and the I of I-Thou is a different I from the I of I-It \parencite{buber1970}. Husserl's Fifth Meditation supplies the transcendental setting: the other is given originally in a paired, analogizing apperception, never as a completed object, so that the other's interiority is constitutively beyond full presentation \parencite{husserl1960}. Sartre's analysis of the look shows what one-directional regard does: under the gaze that I cannot join, I am congealed into an object, my possibilities alienated \parencite{sartre1956}. The phenomenological tradition thus converges on the structure the formal definition encodes: presence is mutual, or it is something else wearing presence's clothes.

The Lacanian layer adds the final, indispensable qualification. Full transparency between lovers is not the ideal; it is the imaginary's characteristic lure. The gaze of the other is never returnable from the place it is sent \parencite{lacan1998}, the other's desire is never exhausted by my image of it, and an intimacy that promises completed mutual knowledge promises a m\'econnaissance. The ideal that survives the qualification is the one the layered account prepared:

\begin{defn}[Presence with acknowledged gap]\label{p05:def:presence}
The other is present, in the constitutive sense, when attentional presence (the standing availability of joint attention) is sustained together with the acknowledgment that the other exceeds every image and every alignment: presence that neither withdraws before the gap nor pretends to have closed it.
\end{defn}

This is the paper's account of what Weil glimpsed in calling attention the rarest and purest form of generosity \parencite{weil1951}: rare, because it gives what cannot be delegated, stored, or feigned without detection; pure, because what it gives is not a holding but a condition, the other's continued formation as a subject. And it is the third paper's conclusion arrived at from the other side: to keep open the room in which she may always answer is exactly to sustain attentional presence across the acknowledged gap.

\subsection{Presence without reciprocity: care and the held-open triangle}
\label{p05:sec:care}

Definition \ref{p05:def:ja} requires mutuality, and the hardest scenes of intimate life are those in which mutuality fails not by lapse but by incapacity: the partner in deep depression, in dementia, in grave illness, who cannot, for a season or forever, occupy the second vertex. If the triangle is the mechanism of formation, does the mechanism simply stop here? Any account that says yes has misdescribed the most serious form of love there is, and two literatures stand ready to press the point. Noddings's analysis of caring places at its center \emph{engrossment}: the one-caring's receptive, non-selective attention to the cared-for, a feeling-with that does not first require symmetrical return \parencite{noddings1984}. Her account is, in the present vocabulary, a description of attention sustained where clause (ii) cannot be answered. Kittay's dependency critique generalizes the charge: theories of justice built for symmetrical, cooperating adults render invisible the dependency relations on which every society rests, and a framework whose elementary structure demands reciprocity risks repeating that error at the dyadic scale \parencite{kittay1999}.

The framework's answer is to distinguish the triangle's \emph{actuality} from its \emph{single-handed maintenance}. When the other cannot attend, the caregiver can still hold the triangle open: continuing to address rather than merely to handle, placing objects before a shared gaze that may not arrive, attending to the other's \emph{possible} attending. Formally, $\mathrm{JA}(a,b,x)$ fails; what persists is the standing offer of its conditions, sustained unilaterally. This is exactly the structure Definition \ref{p05:def:presence} prepared: attentional presence was defined as the standing \emph{availability} of joint attention, not its achievement, and the care case shows the definition's point. The double-video finding illuminates the practice from the other side \parencite{murray1985}: what the diminished partner can still receive, often startlingly late into incapacity, is live responsiveness, contingency tuned to whatever contingency remains; and what the caregiver gives, in tuning to it, is Weil's attention at its limit, generosity with no expectation of return, the purest case of the ur-gift of §\ref{p05:sec:urgift}.

Two consequences follow, and both strengthen theses already on the table. First, the care case is the deepest vindication of the No-Claim Thesis (Proposition \ref{p05:prop:noclaim}): it is precisely because attention can be the object of no claim that the caregiver's attending retains the character of gift rather than the discharge of a debt; an enforceable claim-structure would convert the most generous human practice into specific performance, and Kittay's account of the exploitation of dependency workers shows what societies do once they treat such labor as owed and therefore costless \parencite{kittay1999}. Second, Criterion \ref{p05:crit:attention} must be read against capacity: structural absence is an injustice as a pattern \emph{among the capable}; incapacity within the dyad is not injustice but tragedy. The injustice adjacent to the care case lies one circle out, in the arrangements that leave the caregiver herself unattended. Kittay's principle of \emph{doulia}, that those who care for dependents must themselves be supported in nested relations of care, is in the present vocabulary a demand that the caregiver's own triangles not be starved by the society that benefits from her holding one open alone; its proper category is the constitutive injustice of Section \ref{p05:sec:constitutive}, with society as addressee.

% =====================================================================
\section{The Political Economy of Attention}
\label{p05:sec:economy}
% =====================================================================

\subsection{Scarcity, capture, and the industrial organization of attention}

The economic analysis of attention begins from Simon's observation that in an information-rich world the scarce resource is not information but what information consumes: the attention of its recipients \parencite{simon1971}. On that scarcity an industry was built. Wu's history traces the attention merchant's model from the penny press to the smartphone: capture attention at scale, resell it to advertisers, and refine the techniques of capture into a discipline \parencite{wu2016}. Zuboff describes the data-extractive stage of the model, in which behavioral surplus is harvested to predict and shape future attention \parencite{zuboff2019}; Crary, its temporal ambition, the colonization of every formerly fallow hour \parencite{crary2013}. Stiegler gives the diagnosis its depth: the systematic industrial capture of attention is a \emph{psychopower}, operating on the formation of attention itself, and since attention formation is the formation of persons, its capture is a matter of generational care, not consumer preference \parencite{stiegler2010}. Citton draws the ecological conclusion: attention is not only a scarce resource to allocate but a medium in which subjects and collectives live, so that its degradation is environmental, not merely distributive \parencite{citton2017}.

The present framework sharpens this line of critique with the constitution thesis. If Proposition \ref{p05:prop:mechanism} is right, attention is not one resource among others, downstream of the person who spends it. It is the material of the mechanism by which persons are formed and maintained. An industry that extracts attention therefore does not merely compete with intimate life for a scarce input, as television competed with conversation for evening hours. It competes \emph{at the level of constitution}: what it removes from the household is the very process by which the household's members sustain one another as subjects. Hochschild showed that capitalism could commodify the performance of feeling \parencite{hochschild1983}; Fraser, that capital systematically free-rides on and depletes the reproductive relations it presupposes \parencite{fraser2016}. The attention economy is the limit case of the tendency Fraser names: it depletes not the labor of social reproduction but its mechanism.

\subsection{The parasite thesis: pseudo-shared objects}
\label{p05:sec:parasite}

The critique can be made structurally exact, and the exactness is the point of having a formal definition. Consider what an algorithmic feed offers, measured against Definition \ref{p05:def:ja}.

\begin{defn}[Pseudo-shared object]\label{p05:def:pseudo}
An object $x$ is a pseudo-shared object for parties $a$ and $b$ when it solicits from each the phenomenology of shared attention while the conditions of Definition \ref{p05:def:ja} fail: in particular, (a) mutuality failure: $\mathrm{Att}(a,x)$ and $\mathrm{Att}(b,x)$ hold, but neither party attends to the other's attending, and the object itself absorbs the checking glance that clause (ii) requires; or (b) common-object failure: there is no single $x$ at all, but personalized variants $x_a \neq x_b$, indexically private to each party, presented under the guise of a common world.
\end{defn}

Failure (b) deserves emphasis because it is historically novel. Broadcast media produced genuinely common objects: the evening news was the same news, and a couple could form a triangle over it. The personalized feed cannot be triangulated even in principle, because there is nothing at the third vertex that both parties confront: each receives a stream optimized to that profile of capture, and the streams diverge by design. The deepest structural fact about the attention economy, on the present analysis, is not that it distracts but that it \emph{de-commonizes the world}: it dissolves the class of objects over which triangles can form, while supplying in their place objects engineered to absorb, without returning, the attention that triangles are made of.

The feed also counterfeits the second vertex. It is responsive, endlessly and instantly; it adapts to my attention with a contingency more reliable than any human partner's. But its responsiveness is the responsiveness of the replayed mother in the double-video experiment, perfected: contingent in form, absent in fact \parencite{murray1985}. There is no one whose attending I could attend to; the openness of clause (iii) has no bearer. The feed is, in the strict sense of the layered account, an industrial object lodged at the imaginary layer: a mirror that flatters and captivates, incapable of the triadic layer, and therefore incapable of forming subjects, only of holding them. Millions watching the same viral clip are not in joint attention with one another; the structure is a star, every spoke private, with capture at the hub.

\begin{prop}[Parasite Thesis]\label{p05:prop:parasite}
The attention economy is parasitic on the triadic mechanism of subject formation: it captures the attentional material of the triangle by simulating the triangle's vertices (the responsive other, the common object) while satisfying none of the conditions of joint attention. Its extraction is therefore constitutive, not merely distributive: what is removed is not time or enjoyment but participation in the mechanism by which intimate subjects are formed and sustained.
\end{prop}

The familiar tableau, two partners at one table, each in a separate feed, is on this analysis not a lapse of etiquette. It is a scene in which the household's triadic structure has been disassembled and its components are being metabolized elsewhere: physical co-presence without attentional presence, simultaneous absorption without a common object, the form of togetherness emptied of its mechanism.

\subsection{Attention as the ur-gift: value production and circulation}
\label{p05:sec:urgift}

Paper IV analyzed the intimate gift through a value-circulation framework in which the meaning function $\mu$ proved non-factorizable: the meaning of a gift cannot be decomposed into contributions assignable to each party severally, because it takes its value from its position in the relation. The present analysis supplies what that result lacked, a genetic explanation. Meaning is non-factorizable because the original scene of meaning production is triadic: reference itself is assembled inside joint attention (§\ref{p05:sec:triadic}), so anything that means does so by its position in a triangle, and a triangle is precisely a structure with no bilateral decomposition: clause (iii) of Definition \ref{p05:def:ja} is a property of the between.

Attention is, moreover, not merely one more good that circulates in the gift economy; it is the \emph{ur-gift}, the gift presupposed by every other. Mauss's gifts oblige because something of the giver travels with the thing \parencite{mauss1990}; the present framework can say what travels: the attention congealed in the choosing, the timing, the fit, all of it legible to the recipient only inside the triangle, where the gift is jointly attended as an expression of attending. Three properties qualify attention as the limit case of the gift, and each is a clause of Definition \ref{p05:def:ja} read economically. It is non-delegable (clause (ii) requires \emph{her} attending; a hired proxy changes the object). It is non-storable: attention exists only in its giving, and stockpiled attention is a category error. And it is non-settleable: because its value is positional, no transfer can clear the balance, which is the fourth paper's anti-settlement thesis at its root. Weil's quality-quantity distinction completes the economic picture: attention in the constitutive sense is not an effortful quantity but an orientation, a ``negative effort,'' available only as a whole \parencite{weil1951, weil2002}. What can be measured in minutes is screen time; what forms subjects has no unit.

This yields the framework's verdict on its own title. An ``economics of attention'' is legitimate as the study of attention's scarcity, capture, and circulation; it becomes the very pathology it studies the moment it prices what it studies. The generative-justice tradition supplies the alternative grammar: value generated in unalienated circulation, returning to its generators rather than extracted from them \parencite{eglash2016}. The triadic analysis gives that tradition its smallest exhibit: the intimate triangle is an unalienated attention circuit, the elementary cell of generative value, and the attention economy is its enclosure.

\subsection{The competition for attention: three arenas and a rigged contest}
\label{p05:sec:competition}

The critique so far has faced outward, toward the extractive industry. But competition for attention is not only an external assault; it is the standing internal condition of any finite attender, and a framework that forbade accounting (§\ref{p05:sec:urgift}) while admitting scarcity (§\ref{p05:sec:economy}) owes an explanation of how the two coexist. Three arenas must be distinguished, because the normative situation differs in each.

The first arena is competition \emph{between domains}: work against partner, parents against partner, one's own projects against the shared world. Here Simon's scarcity holds with full force inside the dyad, and the resolution is a two-tier economy already latent in §\ref{p05:sec:imperfect}. The \emph{clearing}, the guarded evening, the protected morning, the device put away, is countable, schedulable, and therefore a legitimate object of distribution: partners can argue justly about clearings, trade them, and owe them. The \emph{attending} that may or may not fill a clearing is not countable and can be the object of no allocation. Distributive justice operates on the first tier; constitutive normativity on the second; and the apparent conflict between scarcity and the anti-accounting stance dissolves, because what is scarce and allocable (cleared time) and what is constitutive and unallocable (attention) were never the same thing. A partner who delivers every negotiated clearing and attends in none of them wrongs the other at the second tier while keeping perfect accounts at the first, which is exactly the phenomenology of certain dutiful, dead marriages.

The second arena is competition \emph{between triangles}. Every subject is a vertex of many: partner, parents, friendships, vocation, and the framework entails that this plurality is not a regrettable dilution but a condition of health, since each triangle forms its subjects and a subject formed in one structure only would be that structure's artifact. It follows that the intimate triangle holds no rightful monopoly. The demand for monopoly has a name, jealousy, and the framework gives its grammar precisely: jealousy is the claim-right to attention (Proposition \ref{p05:prop:noclaim}) reasserted in the register of passion, and its enforcement would be domination in Pettit's sense \parencite{pettit1997}, a standing arbitrary power over the faculty by which the other is a subject at all. The healthy picture is instead cellular: a foam of triangles sharing vertices, attention circulating among cells, each cell sustained by, and contributing to, the formation of subjects it shares with others.\footnote{The cellular picture is developed formally in the author's companion programme on value circulation in relational structures, where the persistence conditions of such cells are treated dynamically.}

The third arena is competition \emph{between the human triangle and its industrial counterfeit}, and here the point missed by Proposition \ref{p05:prop:parasite} taken alone is that the contest is rigged. The human second vertex is metabolic: she tires, needs sleep, has moods, requires that attention be returned for her own attending to be sustained. The feed is a pump: available at every hour, asking nothing, never depleted, and engineered, by iterated optimization against human salience responses, into a supernormal stimulus that no evolved partner can match on the stimulus dimension. This is not competition between like rivals for a common resource; it is a reciprocal circuit competing against a one-way extractor, a metabolism against a siphon. Ordinary liberal vocabulary, consumer choice between alternatives, presupposes symmetrical competitors and so cannot even state the unfairness. The framework can: one competitor must sustain the triadic conditions to offer anything at all, while the other wins precisely by voiding them; the contest is structured so that the constitutive option bears costs the parasitic option externalizes. The unfairness of the contest, and not only the harm of its outcome, is a ground for the third-party duties of §\ref{p05:sec:imperfect}, and for treating the design of capture as answerable conduct rather than neutral supply.

\subsection{The regenerative circuit: a criterion of relational sustainability}
\label{p05:sec:regeneration}

The competition analysis raises the dynamical question the paper has so far suppressed: where does the capacity to attend come from, and what replenishes it? Collins's theory of interaction ritual chains supplies the mechanism \parencite{collins2004}. Successful interaction rituals, defined by mutual focus of attention and shared affect, generate entrainment, and their output is what Collins calls emotional energy: a renewed capacity and motivation for further interaction, which carries forward into the next encounter and makes chains of ritual self-sustaining. Read into the present framework, the finding is exact: \emph{joint attention is regenerative}. Attending together, when it succeeds, produces in each party more capacity for attending than it consumed; the triangle is not only the site where attention is spent but the scene where it is renewed. Extractive circuits invert the sign: attention spent into a feed returns no energy to any relation; what it returns is compulsion, the motivation to repeat the capture, which mimics energy as the pseudo-shared object mimics the common world.

\begin{crit}[Attentional sustainability]\label{p05:crit:sustain}
An intimate relation is attentionally sustainable iff its circuit is net regenerative: the practice of joint attention within it restores, in each party, the capacity and motivation for further attending at least as fast as the relation consumes them, after the demands of all the parties' other triangles and of the extractive environment are paid.
\end{crit}

The criterion unifies three diagnoses the paper has already made. The caregiver of §\ref{p05:sec:care} runs locally net-negative by necessity, holding the triangle open is expenditure without entrainment, which is why Kittay's doulia is not charity but thermodynamics: the caregiver's depletion must be made good from outside the dyad or the held-open triangle eventually falls with its holder. The asymmetric maintainer of §\ref{p05:sec:gender} is the slow-motion case: a relation in which one party chronically produces the conditions the other consumes is unsustainable in exactly the criterion's sense, and its eventual exhaustion is predictable, not mysterious. And the household colonized by feeds (§\ref{p05:sec:parasite}) runs net-negative as a whole, its attentional yield metabolized elsewhere, so that partners may truthfully report having no energy left for one another: the report is not a metaphor but a balance. The criterion is also where the generative-justice tradition acquires its dynamics at the smallest scale \parencite{eglash2016}: value returning to its generators is, in attentional terms, the condition of the cell's persistence, and extraction is not merely unjust but, for the cell, fatal.

% =====================================================================
\section{The Symbolic Order as the Distributor of Attention}
\label{p05:sec:symbolicdist}
% =====================================================================

\subsection{The virtual fourth vertex}
\label{p05:sec:fourth}

Section \ref{p05:sec:symbolic} placed the symbolic mechanism downstream of the triadic in ontogeny; §\ref{p05:sec:transfer} conceded that adult joint attention shares its objects under descriptions. The concession has a consequence that must now be drawn in full, because it reorganizes the political economy just given. The symbolic order is, before it is anything else, an \emph{attention-allocation device}. Salience is socially pre-distributed: a culture's codes settle what is notable and what invisible, which dates demand observance, which features of a face, a house, a career are to be seen at all, who must keep what in mind for whom. No adult chooses against a neutral field of candidates for notice; the field arrives pre-weighted. Lacan's thesis that human desire is the desire of the Other has, in this register, a direct translation: attention follows desire, desire is mediated by the Other, therefore my attention is, in its default settings, the Other's attention forwarded through me. Althusser's hail makes the same point at the level of the event: ``Hey, you there!'' captures attention before it conveys any content; interpellation is attention-capture with a name attached, and the subject's turn toward the call is the elementary act of socially allocated notice \parencite{althusser1971}.

\begin{prop}[Fourth Vertex Thesis]\label{p05:prop:fourth}
Every adult triangle contains a virtual fourth vertex: the symbolic order, which pre-frames the candidate set of shared objects and pre-weights their claims to notice. Adult joint attention is never bare; it is always exercised within an allocation of salience that the parties did not author.
\end{prop}

The thesis is descriptive, not defeatist; what can be done about the allocation is the business of §\ref{p05:sec:autonomy}. Its immediate yield is explanatory. First, it gives §\ref{p05:sec:gender} its etiology: the scripts that assign anticipation and monitoring to women are the fourth vertex at work, a socially legislated distribution of who must hold what in mind, so that the measured asymmetry \parencite{daminger2019} is not an aggregate of individual bargains but the household-level execution of a code. Second, it explains why the couple's common world feels found rather than made: the anniversary, the milestone, the photogenic view present themselves as intrinsically notable because their notability was legislated elsewhere and earlier. Third, it identifies the true depth of the attention economy, pursued in the next subsection: an industry that can write to the fourth vertex does not need to win any particular contest for notice, because it sets the weights by which all contests are scored.

\subsection{Attentional alienation}
\label{p05:sec:alienation}

Jaeggi's reconstruction of alienation provides the concept the framework now needs. Alienation, on her account, is a relation of relationlessness: not the absence of relation but a deficient one, in which one's life, roles, and desires are executed without being \emph{appropriated}, made one's own in the living of them; the criterion of the unalienated life is the capacity for such appropriation \parencite{jaeggi2014}. Transposed to the present subject matter:

\begin{defn}[Attentional alienation]\label{p05:def:alienation}
A party's attention is alienated when its pattern of salience follows a script she has had no part in authoring and cannot recognize herself in: attending that is executed but not appropriated, such that what she notices, tracks, and keeps in mind is, in Jaeggi's sense, hers without being her own.
\end{defn}

The definition would matter little if attention were one activity among others. Under Proposition \ref{p05:prop:mechanism} it matters totally: subjects are formed by what they jointly attend to, so alienated attention entails alienated \emph{formation}. Alienation thereby ceases to be a state that befalls finished subjects, the classical picture in which a completed worker confronts his estranged product, and becomes a pathology of the forming process itself: the subject assembled, day by day, out of notice she cannot own. This is the subjective face of constitutive injustice (Section \ref{p05:sec:constitutive}): where that category describes the arrangement, this concept describes what it is like to be formed inside one, the faint, pervasive sense, familiar from every description of feed-saturated life, of having spent one's day in attendings that belong to no one.

The concept also locates the attention economy's deepest operation, at a depth the parasite thesis does not reach. Capturing instances of attention is foraging; the strategic prize is the \emph{code of attendability} itself: the rewriting of what counts as worth seeing. Paper IV's Baudrillardian analysis showed sign-value to be legislated by position in a differential code \parencite{baudrillard1981}; the platform economy is, in these terms, a private legislature of salience, its trend mechanics and virality metrics functioning as continuous amendments to the code, ratified by no one, binding on the default attention of billions. Stiegler's psychopower names the regime \parencite{stiegler2010}; the present analysis supplies its constitutional location: power exercised at the fourth vertex, upstream of every triangle, legislating the weights under which all common worlds are composed. Enclosure of the attentional commons is therefore the exact phrase: not the theft of attention-instances but the privatization of the salience code through which attention is allocated at all.

\subsection{Attentional autonomy: re-authoring the script}
\label{p05:sec:autonomy}

If the fourth vertex always allocates, is autonomy a consoling fiction? The structure of the answer is the one Paper III built for freedom under the symbolic, and the isomorphism is exact. There, autonomy could not mean exemption from the symbolic order, since the very terms of any vow are socially given; it meant co-legislation within it, the joint authorship of the law one lives under. Here, autonomy cannot mean attention exercised free of any code, a view from nowhere that the layered account of §\ref{p05:sec:lacantomasello} has already exposed as the imaginary's fantasy; it means the recurrent, joint re-authorship of the allocation one attends under. Butler supplies the mechanism's possibility: a code persists only through reiteration, and reiteration is the site of inflection \parencite{butler1997}; couples demonstrably develop attentional idiolects, private canons of the notable, household liturgies, names and dates and objects weighted by no public code, which are precisely deviations written into the script by its local executors.

The device of re-authorship is one the paper already possesses. The reflexive triangle (Definition \ref{p05:def:reflexive}) was introduced as joint attention to the relation; its deepest application is joint attention to \emph{the couple's own pattern of salience}: what do we keep noticing, what have we stopped seeing, whose calendar writes our evenings, which of our observances are ours and which merely arrived. In such a review the script is converted from frame into object, from the invisible allocator of attention into a thing attended; and what is jointly attended can be jointly amended.

\begin{defn}[Attentional autonomy]\label{p05:def:autonomy}
A couple is attentionally autonomous to the degree that the allocation of salience governing their joint attention is recurrently placed under their own reflexive triangle, and survives, or is revised by, that review: autonomy not as code-free attending, which does not exist, but as standing joint legislative power over the code's local form.
\end{defn}

The series closes a circle here. Paper III argued that the normativity of self-legislation in intimate life requires co-authorship of the law; this section identifies the most pervasive law a couple lives under, the law of salience, and finds the same result: legislated unilaterally, whether by one partner, by the inherited code, or by a platform writing to the fourth vertex, it alienates; legislated jointly, through the reflexive practice, it is the couple's freedom in its daily, hourly form. Attentional autonomy is self-legislation continued by attentional means.

% =====================================================================
\section{The Rights Structure of Attention}
\label{p05:sec:rights}
% =====================================================================

\subsection{Against a claim-right to attention}

If attention is constitutive, the partner deprived of it suffers a wrong of the deepest kind, and the jurisprudential reflex is to reach for a right: does $a$ hold, against her partner $b$, a claim-right to attention, with its correlative duty on $b$ to attend? Hohfeld's table gives the question its exact form \parencite{hohfeld1919}, and the answer defended here is no, on three independent grounds.

The first is conceptual impossibility of performance. A claim-right's correlative duty must be dischargeable, and what discharges it must be the thing claimed. But attention in the constitutive sense cannot be produced on demand: compelled attending delivers the posture of clause (i) while voiding clause (iii), since attention rendered under claim is monitoring oneself rendering it, and the openness of the triangle collapses into the performance of openness. What can be compelled is attention's behavioral shell: presence in the room, the face turned, the phone down. The shell is not nothing, and §\ref{p05:sec:imperfect} will return to it, but a right to the shell is not a right to attention, and calling it one obscures what was wanted.

The second ground is the isomorphism with Paper III's central result. A vow, that paper argued, cannot be unilaterally legislated, because its key terms are held in common; a fortiori, one party cannot by unilateral claim fix what shall count as adequate attention from the other. A claim-right to attention would install exactly the unilateral legislature the third paper dismantled, with the claimant as sovereign over the meaning of ``enough.''

The third ground is republican. Pettit's analysis of domination, subjection to another's power of arbitrary interference, applies symmetrically here \parencite{pettit1997}: a legally or morally enforceable claim on the other's inner orientation would constitute a standing power over the very faculty by which the other is a subject, and would dominate precisely what it meant to cherish. The beloved who attends under enforceable claim is not present; she is occupied.

\begin{prop}[No-Claim Thesis]\label{p05:prop:noclaim}
Within the Hohfeldian table, constitutive attention can be the object of no claim-right. The intimate parties hold, with respect to each other's attention, privileges (each is at liberty to offer or seek attention) and powers (each can, by invitation, vow, or ritual, alter the normative landscape of their attending), but no claims; correlatively, each stands under no duty enforceable by the other's unilateral demand.
\end{prop}

\subsection{Imperfect duty and the structural criterion of attentional justice}
\label{p05:sec:imperfect}

The No-Claim Thesis does not leave attention normatively naked; it relocates the normativity. Kant's category of imperfect duty fits with unusual precision \parencite{kant1996}: duties of love, in his architecture, are of wide obligation, prescribing maxims and orientations rather than enumerable acts, leaving latitude in execution, and admitting of no external compulsion. The duty to attend in intimate life is of this kind: a standing orientation toward attentional presence whose particular occasions cannot be specified in advance, owed not because claimed but because one has, by the relation itself, undertaken the maintenance of its triangle.

The duty register is, however, only half the ethics, and the lesser half. Murdoch, inheriting Weil explicitly, made attention the central concept of morals: the characteristic work of the moral agent is a ``just and loving gaze'' directed upon an individual reality, and moral progress is the slow defeat of fantasy by attention, learning to see the other as she is rather than as one's need composes her \parencite{murdoch1970}. On her account attention is not preparation for ethical action but the primary ethical act, and choice, when its moment arrives, is largely the residue of prior seeing: by the time we choose, the quality of our attention has mostly decided what we are able to choose between. For intimate life the consequence is that a partner's attending is itself a moral achievement or failure, prior to and independent of any act it issues in; and the imaginary layer of §\ref{p05:sec:lacantomasello} names exactly the adversary Murdoch names, the self-serving fantasy that composes the other, against which attention is the discipline. The ethics of attention is therefore double-registered: an imperfect duty, in Kant's architecture, to maintain the orientation; a virtue, in Murdoch's, that the orientation, practiced, becomes. The epigraph of this paper is on this reading not an ornament but a thesis.

Justice enters not at the level of acts but of structure. The criterion is built on the concept that did analogous work in Fricker's theory: a wrong identifiable by a structural pattern rather than a culpable episode \parencite{fricker2007}.

\begin{crit}[Attentional injustice]\label{p05:crit:attention}
An intimate relationship exhibits attentional injustice when there is a standing, asymmetric pattern in which one party's concerns, projects, and experiences systematically fail to become objects of joint attention: not the fluctuation of busy seasons, but a structure in which what matters to one partner reliably cannot reach the triangle. The test is structural absence over time, never a quantity of minutes; the wrong is identified by what the relation's attentional pattern renders chronically invisible, not by an accounting of attention paid.
\end{crit}

The criterion deliberately echoes hermeneutical injustice, and the relation can now be stated more strongly than kinship: attentional injustice is constitutively \emph{upstream} of epistemic injustice. Attention is the gate of evidence. What never reaches joint attention is never articulated, since articulation is downstream of the triangle (Vygotsky's law, run in reverse); what is never articulated never becomes testimony; what never becomes testimony never enters the interpretive commons from which collective hermeneutical resources are built. Each of Fricker's two injustices accordingly has an attentional etiology. Testimonial injustice operates, before any credibility judgment, as attention withdrawal: the speaker not looked at as a knower, the gaze withdrawn before the words are weighed, so that the credibility deficit is often executed attentionally without ever being formulated. Hermeneutical injustice begins, generations earlier, in standing patterns of what is never jointly attended: the experiences for which a culture lacks concepts are, with few exceptions, the experiences its triangles were structured never to take as objects, and the fourth vertex of §\ref{p05:sec:fourth} is the mechanism of the structuring. The two wrongs remain distinct: attentional injustice within a dyad can obtain in full where the collective hermeneutical commons is intact, and conversely. But the direction of constitution runs one way, and a theory of epistemic justice that starts at credibility and concepts has started two steps downstream of the wrong. Section \ref{p05:sec:constitutive} gives the wrong its proper genus.

Two corollaries close the section. First, the behavioral shell recovers a modest enforceability: partners can legitimately bind themselves, by the powers Proposition \ref{p05:prop:noclaim} preserves, to structures that protect the triangle's conditions, the guarded evening, the device-free table, the standing review; what is enforced is the clearing, never the attending. Second, the criterion generates third-party duties: if attentional injustice is identified structurally, then designers of attention-capturing systems are answerable for foreseeably producing, at scale, the structural absence the criterion names. The jurisprudence of that answer belongs to the companion paper on AI-mediated intimacy; the ground for it is laid here.

\subsection{The gendered distribution of attentional labor}
\label{p05:sec:gender}

Criterion \ref{p05:crit:attention} is formally neutral between the parties; actual households are not, and a criterion that ignored this would let formal neutrality launder a substantive asymmetry. The empirical record is unambiguous. Hochschild documented the second shift and the emotion work that runs it, disproportionately performed by women \parencite{hochschild1989, hochschild1983}; Daminger's decomposition of cognitive household labor into anticipation, identification, decision, and monitoring found the anticipatory and monitoring components, the labor of \emph{keeping things in mind}, markedly skewed by gender even in couples whose physical task division is near parity \parencite{daminger2019}.

In the present vocabulary this labor has an exact location: it is \emph{triangle maintenance}. Anticipating, monitoring, keeping the household's objects within reach of joint attention, noticing what must be noticed: this is the work of holding the third vertex stocked and the structure standing, and it is invisible for a principled reason: its product is not a thing but a structure, and a structure well maintained presents itself as no work at all. Its etiology was given in §\ref{p05:sec:fourth}: the assignment of this labor is the fourth vertex at work, a socially legislated distribution of who must hold what in mind, executed in households that never voted on it. The asymmetry is therefore double, and the doubling is the diagnosis. The same partner, statistically she, performs more of the attending that maintains the triangle, \emph{and} her own concerns less often become its objects: producer of the mechanism, minority beneficiary of its use. Each half conceals the other, since the partner who curates the couple's common world is thereby positioned as its staff rather than its subject, the one who points rather than the one pointed for.

Criterion \ref{p05:crit:attention} accordingly takes a distributional supplement:

\begin{crit}[Attentional injustice, distributional clause]\label{p05:crit:distribution}
Attentional injustice includes not only the structural absence of one party's concerns from the triangle (Criterion \ref{p05:crit:attention}) but a standing asymmetry in who bears the cost of the triangle's maintenance: an arrangement in which one party chronically produces the conditions of joint attention that the other chiefly consumes.
\end{crit}

Three remarks fix the clause's relations to the rest of the framework. It extends, to the formative structure itself, the third paper's exploitation criterion for emotional labor: what is appropriated here is not surplus feeling but surplus \emph{maintenance}, and Fraser's diagnosis of capital's free-riding on reproduction \parencite{fraser2016} reappears at household scale, one partner free-riding on the reproduction of the relation. It does not breach the No-Claim Thesis: the clause identifies a structural wrong and licenses renegotiation through the powers Proposition \ref{p05:prop:noclaim} preserves, the binding clearings of §\ref{p05:sec:imperfect}; it does not arm anyone with a claim to the other's attending. And it disciplines the care analysis of §\ref{p05:sec:care} against sentimental misuse: the held-open triangle is the glory of care when incapacity compels it, and the costume of exploitation when capacity merely declines, and the difference between the two is exactly what the distributional clause is built to detect.

% =====================================================================
\section{Constitutive Injustice}
\label{p05:sec:constitutive}
% =====================================================================

\subsection{From constitution to justice: two bridges}
\label{p05:sec:bridge}

Before the category is defined, a debt incurred throughout the paper must be paid. The argument has moved from facts of constitution (subjects are formed in triangles) to verdicts of justice (starving the triangle wrongs them), and no fact entails a norm by itself. Air constitutes our breathing, yet foul air is, without more, a misfortune rather than an injustice; constitution alone confers no normative standing. Shklar's line marks what the inference still needs: misfortune is what no arrangement answers for, injustice is what some arrangement does \parencite{shklar1990}. Two bridges carry the argument across, independent of each other and converging on the same category; their independence matters, because they will turn out to ground duties for different addressees.

The first bridge is transcendental, and runs through the practice of justice itself. Every theory of justice, whatever its content, presupposes subjects: holders of holdings, claimants of claims, bearers of the standing to be wronged. The conditions under which subjects are formed and maintained therefore occupy a position no ordinary good can occupy: they are presupposed by the very practice within which goods are distributed and claims adjudicated. One cannot coherently operate a scheme of justice among subjects while remaining indifferent to the mechanism that makes and keeps them subjects, any more than one can value the moves of a game while disclaiming all interest in there being a board. If anything within the practice matters, the preconditions of the practice matter, and not derivatively, as one weighted good among others, but structurally, as what the weighing presupposes. Joined to Shklar's line, the bridge yields the inference the paper has been using: where the degradation of constitutive conditions is the work of an \emph{arrangement}, something designed, maintained, profited from, and revisable, it crosses from misfortune into injustice, because there is now an address to which the complaint can be delivered. This is why the verdicts of §\ref{p05:sec:parasite} are verdicts of justice: the pseudo-shared object is engineered.

The second bridge is fiduciary, and is native to this series. To enter an intimate relation is, on the analysis of §\ref{p05:sec:reflexive}, to build a structure of which clause (iii), mutual openness, is a fact about the between: each party thereby comes to hold, in part, something that is not hers alone and that the other cannot hold without her. The grammar of such holding is trust. Obligations of trusteeship arise wherever one keeps what is another's, or what is common, in one's keeping; and the partners in a triangle are, each toward the other, trustees of a formative mechanism that exists only as jointly held. This grounds the phrase that §\ref{p05:sec:imperfect} used on credit, that the duty to attend is owed ``by the relation itself'': the relation is the trust, and the imperfect duty of attention is its fiduciary content. The bridge is the third paper's result run one level deeper: there, the vow's terms were held in common and so could not be unilaterally legislated; here, the formative structure itself is held in common and so cannot be unilaterally neglected, since neglect by either party degrades what both hold.

The two bridges divide the normative labor cleanly. The fiduciary bridge binds intimates, with the special stringency of trust, and grounds the dyadic applications of the category (Criteria \ref{p05:crit:attention} and \ref{p05:crit:distribution}). The transcendental bridge binds strangers, designers, and institutions, who stand in no trust relation to a given couple but who operate arrangements upon the preconditions of subjecthood at large; it grounds the wider addressees of §\ref{p05:sec:scope}. A wrong reachable by both bridges at once, the spouse who is also the household's attention-extracting designer, is correspondingly compounded. With the bridges in place, the category can be defined without suppressed premises.

\subsection{The category defined}

The wrongs assembled in the preceding sections, the structural starvation of the triangle within a relationship, the industrial enclosure of attention from without, share a form that existing categories of injustice do not capture. Distributive injustice concerns holdings; rights violation concerns enumerated claims; Honneth's recognition wrongs concern the withholding of due regard from a formed subject seeking confirmation \parencite{honneth1995}; Fricker's epistemic injustices concern a person's standing and resources as a knower \parencite{fricker2007}. The wrong in view here lies beneath all four: it touches neither what a subject has, nor what she may demand, nor how she is regarded, nor what she can know, but the conditions under which she is and remains a subject.

\begin{defn}[Constitutive injustice]\label{p05:def:ci}
An arrangement is constitutively unjust toward a person when it systematically degrades her access to, or participation in, the mechanisms of subject formation and maintenance, paradigmatically the triadic mechanism of joint attention, independently of any loss of holdings, violation of enumerated rights, denial of recognition, or epistemic exclusion that may also occur.
\end{defn}

The independence clause carries the weight, so its work is shown by separation cases. A person may be distributively well off, legally unviolated, publicly esteemed, and epistemically credited, while living inside an arrangement, a marriage of parallel feeds, a workplace engineered for capture, that starves every triangle she belongs to: constitutive injustice without the other four. Conversely, a poor and disregarded household may sustain triangles of full intensity. The category is therefore not reducible to its neighbors, though in practice it compounds them, since the subject degraded in formation is thereby weakened in claiming, knowing, and seeking recognition: constitutive injustice is upstream injustice, and its effects launder themselves through the categories downstream.

Adjacency to Honneth deserves one more sentence, since love is his first sphere of recognition and the kinship is close \parencite{honneth1995}. Recognition, in his grammar, is something a subject receives and can be denied; the triangle, in the present grammar, is something a subject inhabits and can be evicted from. Misrecognition insults; constitutive injustice unmakes. The two will usually travel together, but the second is the deeper stratum, and naming it separately is what allows the critique of the attention economy to be stated as a thesis about justice rather than a complaint about manners.

\subsection{Scope and addressees}
\label{p05:sec:scope}

Constitutive injustice, so defined, has three addressees, in widening circles. Within the dyad, it grounds Criterion \ref{p05:crit:attention} and the imperfect duties of §\ref{p05:sec:imperfect}: partners are the first trustees of each other's formation. At the level of design, it indicts architectures that produce pseudo-shared objects at scale (Proposition \ref{p05:prop:parasite}): the wrong done by an attention-extractive system is not adequately described as manipulation of choice or appropriation of data, vocabularies that presuppose a finished subject defending her perimeter; it is the strip-mining of the material from which subjects are made, and its proper legal register is closer to the protection of conditions (environmental, in Citton's sense \parencite{citton2017}) than to the vindication of transactions. At the level of political theory, it names what a society owes its members beyond distribution and rights: the protection of formative structures, of which the intimate triangle is the elementary case and public institutions of common attention, the shared text, the common deliberation, are the civic enlargements.

One boundary is drawn against overreach. Not every solitude, austerity, or chosen withdrawal from relation is constitutive injustice; the category requires systematic degradation of access imposed by an arrangement, not the exercise of a person's own privileges over her attention. A theory that conscripted everyone into triangles would violate Proposition \ref{p05:prop:noclaim} at social scale. The object of justice is the standing availability of the mechanism, not its compelled use.

% =====================================================================
\section{Objections and Replies}
\label{p05:sec:objections}
% =====================================================================

\begin{objection}[Romantic technophobia]
The argument repeats a perennial moral panic. Print, radio, and television were each accused of dissolving intimate life; households absorbed them and formed their triangles around them. The feed will be domesticated in turn.
\end{objection}
\begin{reply}
The reply is in Definition \ref{p05:def:pseudo}(b), and it is structural, not nostalgic. Broadcast media were attention-hungry but \emph{common-object-producing}: the same program reached every receiver, and households did in fact triangulate over it. The personalized feed is the first mass medium whose objects are indexically private by design; there is no common $x$ for the triangle to take. The argument does not predict that domestication is impossible, and §\ref{p05:sec:imperfect}'s enforceable clearings describe its form; it claims that domestication now requires deliberately rebuilding the common object that earlier media supplied for free. The asymmetry is in the technology's structure, not in the critic's mood.
\end{reply}

\begin{objection}[The empirical hostage]
The paper rests normative architecture on a live empirical literature. Shared intentionality, the nine-month revolution, even joint attention's specialness are contested; if the psychology shifts, the ethics collapses.
\end{objection}
\begin{reply}
The dependence is narrower than it appears. The normative sections require three claims: that joint attention as defined exists and is achieved (uncontested); that it is implicated in the formation and maintenance of subjects (supported convergently by the developmental record, by Vygotsky's law, and, independently of all psychology, by Davidson's transcendental argument \parencite{davidson2001}); and that its conditions can structurally fail (Definition \ref{p05:def:pseudo} is an analysis, not an experiment). Disputes internal to the literature, over the uniqueness of human shared intentionality or the mechanism of the nine-month transition, do not touch these. The paper holds the psychology as a load-bearing wall, but the building has a second foundation in the transcendental argument, and the two would have to fail together.
\end{reply}

\begin{objection}[Levinas redux]
The shared object is a thing; Levinas's third is a face. Substituting world for tiers domesticates the very asymmetry that makes the ethical ethical, and licenses a complacent symmetry between lovers.
\end{objection}
\begin{reply}
The substitution was offered as structural, not ethical, and the layered account keeps the registers apart. The shared object performs the third's structural function, instituting comparability, which is all that justice-talk within the dyad requires. The ethical asymmetry of the face is untouched: it lives at a layer the triangle neither generates nor cancels, and Definition \ref{p05:def:presence} writes its trace into the ideal itself, as the acknowledged gap. The position is not that justice replaces the asymmetrical ethical, but that the dyad is large enough to contain both, because its own mechanism gives it three vertices.
\end{reply}

\begin{objection}[The accountant's return]
Criterion \ref{p05:crit:attention} forswears minutes but speaks of patterns over time. A pattern is detected by tracking; tracking is accounting; the framework smuggles back the settlement logic Paper IV expelled.
\end{objection}
\begin{reply}
The criterion identifies a wrong by structural absence: a class of concerns that \emph{cannot reach} the triangle. Detecting that something never appears requires no ledger of what does; one audits the silence, not the traffic. The relevant practice is itself triadic, the standing review in which the relation is placed under joint attention (§\ref{p05:sec:reflexive}) and each party may say what has become invisible. Where Paper IV's settlement logic priced contributions to clear them, the reflexive review prices nothing and clears nothing; it restores objects to the triangle. The difference between accounting and witnessing was drawn in the fourth paper; the criterion stands on the witnessing side.
\end{reply}

\begin{objection}[Symbolic determinism]
The Fourth Vertex Thesis proves too much. If the symbolic order allocates salience before any choice, then attentional autonomy (Definition \ref{p05:def:autonomy}) is a consoling fiction: the couple's reflexive review is itself conducted in inherited concepts, scheduled by inherited calendars, and so merely executes the script at a higher level of recursion.
\end{objection}
\begin{reply}
The objection restates, for attention, the determinist's case against self-legislation, and the reply is the compatibilist structure Paper III already built. Autonomy was never claimed as exemption from the symbolic; it was defined (Definition \ref{p05:def:autonomy}) as joint legislative power over the code's \emph{local form}, exercised from within it, and an inside position is the only position there is. Three facts show the power to be real rather than notional. Codes persist only through reiteration, and reiteration admits inflection \parencite{butler1997}: drift is observable at every scale from slang to liturgy, and someone's reiterations are doing the drifting. Couples observably develop attentional idiolects, private canons of the notable that no public code weights, which are deviations written by local executors and inherited by no one else. And the reflexive review can \emph{fail}: couples attempt it and find their salience pattern unbudged, which would be impossible if the review merely executed the script, since execution cannot fail against itself. What the objection demands, attention legislated from no code at all, is the view from nowhere, the imaginary's fantasy of a subject prior to all formation, and the layered account discharged it in §\ref{p05:sec:lacantomasello}. Freedom under the fourth vertex is the only freedom on offer, and it is enough: the difference between a couple that reviews its law of salience and one that merely lives it is the difference Paper III drew between co-authored and unilateral legislation, here applied to the legislature that sits longest.
\end{reply}

% =====================================================================
\section{Conclusion}
\label{p05:sec:conclusion}
% =====================================================================

The paper has argued one thesis and traced its consequences. The thesis: the subject of intimate life is formed and sustained through the presence of the other and the practice of joint attention; the triadic mechanism is the hinge between the mirror in which the self first appears and the symbolic order in which it takes a name (Propositions \ref{p05:prop:mechanism}, \ref{p05:prop:hinge}). The mechanism transfers to adult intimacy at reduced intensity but unchanged structure, finds its mature form in the reflexive triangle, where the relation itself becomes the shared object (Proposition \ref{p05:prop:transfer}, Definition \ref{p05:def:reflexive}), and shows its limit form in care, where one party holds the triangle open for an other who cannot, for a time, answer (§\ref{p05:sec:care}). The consequences: the attention economy is a parasite on the mechanism, extracting constitutive material through pseudo-shared objects in a contest rigged against the metabolic competitor, and relational sustainability has a criterion, the net-regenerative circuit (Proposition \ref{p05:prop:parasite}, Criterion \ref{p05:crit:sustain}); the symbolic order is the standing distributor of attention, the virtual fourth vertex of every triangle, so that alienation reappears as a pathology of formation itself and autonomy as the couple's joint legislative power over its law of salience (Proposition \ref{p05:prop:fourth}, Definitions \ref{p05:def:alienation}, \ref{p05:def:autonomy}); the normativity of attention within the dyad takes the form not of claim-rights but of imperfect duty and Murdochian virtue, disciplined by a structural criterion of attentional justice, read always with its distributional clause and standing constitutively upstream of epistemic injustice (Proposition \ref{p05:prop:noclaim}, Criteria \ref{p05:crit:attention}, \ref{p05:crit:distribution}); and the theory of justice gains an upstream category, constitutive injustice, the degradation of the conditions of subjecthood itself, carried from fact to norm by two independent bridges, transcendental for strangers and institutions, fiduciary for intimates (Definition \ref{p05:def:ci}, §\ref{p05:sec:bridge}).

Within the series, the paper closes two circles that the earlier papers opened. The vow of Paper III and the gift of Paper IV are now legible as two technologies of the same structure: devices for placing the relation under joint attention, which is why neither could be unilateral and neither could be settled. And the third paper's deepest result returns at the end of this one in generalized form: the most pervasive law a couple lives under is its law of salience, and the alternative to its joint legislation is not lawlessness but alienation, government by a code amended daily, at the fourth vertex, by parties the couple never seated. What the earlier papers presupposed, this paper supplies: an account of the subjects whose vows and gifts they are, formed in the triangle they keep rebuilding, free exactly insofar as they keep its law their own. The next perimeter is already visible. If subjects are formed in triangles, and machines now offer themselves as the second vertex, then the question of what an artificial interlocutor can and cannot occupy in the structure of Definition \ref{p05:def:ja}, clause by clause, is not futurism but the present paper's unfinished business, and the companion work on AI-mediated intimacy is where it falls due.

The oldest description of love's beginning in the Chinese canon is a scene of joint attention under the night sky: \zh{今夕何夕，见此良人}, what night is this, that I see this good person. Not possession, not even promise: seeing, in a night both stand inside. The argument of this paper is that the old scene is exact. To love is to keep a world in common, and to keep attending, together, across the gap that keeps the other forever more than what is seen.

% =====================================================================
\section*{Acknowledgements}
\addcontentsline{toc}{section}{Acknowledgements}
% =====================================================================

This paper is the fifth in a series on the philosophy of intimacy and the theory of justice. The lovers discussed in these pages are kept deliberately anonymous, and every case is presented as a constructed one, so that the argument may stand on its structure rather than on any private history.

My deepest gratitude is owed, as in the papers before it, to the ``forest girl'': pure and natural in spirit, good-hearted, resilient, and wise, carrying within her a quiet sense of mission toward the world. This paper's thesis was learned from her before it was argued: that one becomes most fully oneself in the presence of an other who truly attends. Whatever is true in these pages was first true at a shared table, before a shared window, under a shared moon. The vow renewed in the third paper is renewed here in this paper's own vocabulary: to remain present, to keep the world between us common, and to honor the gap by which she will always exceed what I see; that I will not fail her, to the end of our days.

\begin{center}
{\zh{愿天下有情人，目之所注皆有回应，心之所向皆得见证；同观一月，白首不相负。}}
\end{center}

% =====================================================================
